% !TeX root = ../../main.tex
% =====================================================================
%  Interacción entre aerodinámica y suspensión
%  Responsable: 
%  Última revisión: 
% =====================================================================
\chapter{Interacción entre aerodinámica y suspensión}
\label{ch:interaccion_aero_suspension}

\begin{objetivos}
  \item \redactar{Objetivo de aprendizaje 1 --- concreto y comprobable.}
  \item \redactar{Objetivo 2.}
  \item \redactar{Objetivo 3.}
\end{objetivos}

\fichasesion{90 min}{\redactar{capítulos previos}}{\redactar{material}}

\section{Sensibilidad a altura y cabeceo}
\label{sec:interaccion_aero_suspension:01}
\pendiente{Pendiente de redactar.}

\section{Control de plataforma}
\label{sec:interaccion_aero_suspension:02}
\pendiente{Pendiente de redactar.}

\section{Efecto sobre el balance}
\label{sec:interaccion_aero_suspension:03}
\pendiente{Pendiente de redactar.}

\section{Compromisos de puesta a punto}
\label{sec:interaccion_aero_suspension:04}
\pendiente{Pendiente de redactar.}

% ---------------------------------------------------------------------
%  Cierre del capítulo: los cuatro recuadros son obligatorios.
% ---------------------------------------------------------------------

\begin{caso}
\redactar{Qué hicimos nosotros: decisión tomada, datos de partida, resultado en pista y qué haríamos distinto.}
\end{caso}

\begin{ojo}
\begin{itemize}
  \item \redactar{Error que repetimos cada temporada en este tema.}
\end{itemize}
\end{ojo}

\begin{entregable}
\redactar{Qué produce el lector al terminar: plantilla, hoja de cálculo, checklist o apartado del informe.}
\end{entregable}

\begin{juez}
\begin{enumerate}
  \item \redactar{Pregunta tipo del design event sobre este capítulo.}
\end{enumerate}
\end{juez}
